\documentclass[10pt,a4paper]{amsart}
\usepackage[margin=19mm]{geometry}
\usepackage{amsmath,amssymb,mathtools}
\usepackage[scr=rsfs,cal=euler]{mathalfa}
\usepackage{xfrac}
\usepackage[unicode,hidelinks,bookmarks=false]{hyperref}
\newcommand{\ie}{i.e.\ }
\newcommand{\cf}{cf.\ }
\newcommand{\resp}{respectively\ }
\newcommand{\Subsection}{\subsection}
\providecommand{\nobreakdash}{\nobreak\hbox{-}}
\setlength{\emergencystretch}{2em}
\multlinegap0pt
\usepackage{float}
\usepackage{xcolor}
\newtheorem{lemma}{Lemma}
\title{Preventing large-scale avalanches in the two-dimensional Abelian sandpile model through strategic interventions}
\author{Maike C. de Jongh, Richard J. Boucherie, M.N.M. van Lieshout}
\date{Source: arXiv:2603.24459v1 (CC BY 4.0)}
\begin{document}
\maketitle

\noindent\emph{Source and licence.} Preventing large-scale avalanches in the two-dimensional Abelian sandpile model through strategic interventions. Version arXiv:2603.24459v1 (CC BY 4.0), distributed by arXiv under the Creative Commons Attribution 4.0 International licence, http://creativecommons.org/licenses/by/4.0/, as recorded in the arXiv licence metadata for this exact version and shown on the abstract page https://arxiv.org/abs/2603.24459v1. Full licence notice: \texttt{licenses/arxiv-2603.24459-CC-BY-4.0.txt}.\par\smallskip

\noindent\textit{Editorial scope.} Counted unit: the article's lemma exp\_av\_size\_rec\_thm (source0.tex lines 267--275). The article's complete original proof of it is delivered with it (source0.tex lines 276--308, one \texttt{proof} environment of 33 lines). No mathematics is written, repaired or re-derived: the statement and the proof are the article's own lines, in the article's own notation, and no retained line is substituted or re-flowed.  Retained uncounted as the source-local context the proof needs: lines 222--224.  Nothing was omitted from any retained span, so the package carries no omission record.  No retained line contains a citation, so the wrapper carries no bibliography and no bibliography entry.  No retained line contains a figure environment, an image-inclusion command or drawing code, so no image is delivered and the package declares no asset dependency.  Editorial formatting only, in addition: an AMS \texttt{amsart} preamble that restates the article's own declarations and macros used by the retained lines, namely the environments \texttt{lemma} on separate counters, together with the standard packages those lines need.  The retained environments print in this document as Lemma 1 in order of appearance, whereas the article prints the same items with the numbering of its own full document.  The complete native source of the article as distributed in the arXiv e-print for this exact version is bundled unchanged as \texttt{sources/197147/source0.tex}.
\begin{equation}\label{sum_of_waves}
    x(\eta, v_i) = \sum\limits_{j = 1}^{n(\eta, v_i)} w_{ij}(\eta).
\end{equation}

\begin{lemma}\label{exp_av_size_rec_thm}
  Let $\eta \in \Omega$ denote a stable sandpile configuration and let $A \subseteq V$ denote a generator in $\eta$. Let $Y$ denote the random vertex at which the next sand grain lands. The expected value of the size of the next avalanche corresponding to $A$ satisfies
  \begin{equation}\label{exp_av_size_rec_eq}
      \mathbb{E}[X(\eta)|Y \in A] = \begin{cases}
          w_A(\eta) + \sum\limits_{j=1}^{\ell} \dfrac{|A_j|}{|A|} \mathbb{E}[X(\eta^A)|Y \in A_j], &\text{if } \{v \in A|\eta^A(v) = 3\} \neq \emptyset, \\
          w_A(\eta), &\text{otherwise}. 
      \end{cases}
  \end{equation}
\end{lemma}

\begin{proof}
    Let $\tilde{H}$ denote the subset of vertices in $A$ that, when hit by a new sand grain, generate an avalanche consisting of a single wave. Let $H' = A \setminus \tilde{H}$. If $H' \neq \emptyset$, let $H_1, \ldots, H_{\ell}$ denote its connected components. By conditioning on $Y$, we obtain
    \begin{align*}
        \mathbb{E}[X(\eta)|Y \in A] &= \dfrac{|\tilde{H}|}{|A|} \mathbb{E}[X(\eta)|Y \in \tilde{H}]  + \sum\limits_{j=1}^{\ell} \dfrac{|H_j|}{|A|} \mathbb{E}[X(\eta)|Y \in H_j] \\
        &= \dfrac{|\tilde{H}|}{|A|} w_A(\eta) + \sum\limits_{j=1}^{\ell} \dfrac{|H_j|}{|A|}\mathbb{E}[X(\eta)|Y \in H_j],
    \end{align*}
    if $H' \neq \emptyset$ and
    \begin{equation*}
        \mathbb{E}[X(\eta)|Y \in A] = w_A(\eta),
    \end{equation*}
    otherwise. Hence, to establish expression (\ref{exp_av_size_rec_eq}), it suffices to show that $\tilde{H} = \tilde{A}$, $H_j = A_j$ for all $j = 1, \ldots, \ell$ and
    \begin{equation}\label{exp_av_size_rec_help}
        \mathbb{E}[X(\eta)|Y \in A_j] = w_A(\eta) + \mathbb{E}[X(\eta^A)|Y \in A_j].
    \end{equation}
    We start by showing that $\tilde{H} = \tilde{A}$. First, consider a vertex $v_i \in \tilde{H}$. Let $(x_1, \ldots, x_k)$ be a permutation of $W_A(\eta)$. Since an avalanche caused by dropping a sand grain at $v_i$ consists of only one wave, we have $\tau_{x_k} \cdots \tau_{x_1} \alpha_i \eta(v_i) \leq 3$. Since $\tau_{x_k} \cdots \tau_{x_1} \alpha_i \eta = \alpha_i \tau_{x_k} \cdots \tau_{x_1} \eta$, this implies that $\tau_{x_k} \cdots \tau_{x_1} \eta(v_i) = \eta^A(v_i) < 3$. Hence, $v_i \in \tilde{A}$. Now, suppose that $v_i \in \tilde{A}$. This implies that $\eta^A(v_i) < 3$. Therefore, we have $\tau_{x_k} \cdots \tau_{x_1} \alpha_i \eta(v_i) \leq 3$ for each permutation $(x_1, \ldots, x_k)$ of $W_A(\eta)$. It follows that an avalanche generated by dropping a sand grain at vertex $v_i$ consists of only a single wave, and thus $v_i \in \tilde{H}$. Thus, we have $\tilde{H} = \tilde{A}$.

    We proceed to prove that $H_j = A_j$ for all $j = 1, \ldots, \ell$. Let $A' = \bigcup_{j = 1}^{\ell} A_j$. Note that it suffices to show that $H' = A'$. First, consider $v_i \in H'$. Let $(x_1, \ldots, x_k)$ be a permutation of $W_A(\eta)$. Since the avalanche that is generated after a new sand grain drops at $v_i$ can be decomposed into more than one wave, we have $\tau_{x_k} \cdots \tau_{x_1} \alpha_i \eta(v_i) > 3$. Since each vertex in $W_A$ topples only once, this implies $\tau_{x_k} \cdots \tau_{x_1} \alpha_i \eta (v_i) = 4$. Hence, $\tau_{x_k} \cdots \tau_{x_1} \eta(v_i) = \eta^A(v_i) = 3$. It follows that $v_i \in A'$. Now, suppose that $v_i \in A'$. Hence, $\eta^A(v_i) = 3$. This implies that $\tau_{x_k} \cdots \tau_{x_1} \alpha_i \eta (v_i) = 4$ for each permutation $(x_1, \ldots, x_k)$ of $W_A(\eta)$. Therefore, if a new sand grain lands at $v_i$, the avalanche that is generated consists of multiple waves. Hence, $v_i \in H'$. Thus, we obtain $H' = A'$. 

    Finally, we show the validity of expression (\ref{exp_av_size_rec_help}). Through further conditioning on $Y$, we obtain
    \begin{equation*}
        \mathbb{E}[X(\eta)|Y \in A_j] =  \dfrac{1}{|A_j|}\sum\limits_{v_i \in A_j} x(\eta, v_i) . 
    \end{equation*}
    Suppose that the avalanche caused by dropping a sand grain at vertex $v_i$ consists of $k_i$ waves. Invoking expression (\ref{sum_of_waves}), we obtain
    \begin{equation*}
        \mathbb{E}[X(\eta)|Y \in A_j] = \dfrac{1}{|A_j|} \sum\limits_{v_i \in A_j} \left[w_A(\eta) + \sum\limits_{n=2}^{k_i} w_{in}(\eta) \right].
    \end{equation*}
    Let $(x_1, \ldots, x_k)$ denote a permutation of $W_A(\eta)$. Since $\tau_{x_k} \cdots \tau_{x_1} \alpha_i \eta = \alpha_i \tau_{x_k} \cdots \tau_{x_1} \eta = \alpha_i \eta^A$, it follows that $w_{in}(\eta) = w_{i(n-1)}(\eta^A)$ for all $n = 2, \ldots, k_i$. Hence, we obtain
    \begin{align*}
        \mathbb{E}[X(\eta)|Y \in A_j] &= \dfrac{1}{|A_j|} \sum\limits_{v_i \in A_j} \left[w_A(\eta) + \sum\limits_{n=1}^{k_i-1}w_{in}(\eta^A)\right] = w_A(\eta) + \dfrac{1}{|A_j|} \sum\limits_{v_i \in A_j} \sum\limits_{n=1}^{k_i-1} w_{in}(\eta^A) \\
        &= w_A(\eta) + \dfrac{1}{|A_j|} \sum\limits_{v_i \in A_j} x(\eta^A, v_i) = w_A(\eta) + \mathbb{E}[X(\eta^A)|Y \in A_j].
    \end{align*}
    This concludes the proof.
\end{proof}

\end{document}
